\documentclass[11pt]{article}
\usepackage[margin=1in]{geometry}
\usepackage{amsmath,amssymb,amsthm,booktabs,hyperref}

\newtheorem{definition}{Definition}[section]
\newtheorem{axiom}[definition]{Axiom}
\newtheorem{assumption}[definition]{Assumption}
\newtheorem{conjecture}[definition]{Conjecture}
\newtheorem{lemma}[definition]{Lemma}
\newtheorem{proposition}[definition]{Proposition}
\newtheorem{theorem}[definition]{Theorem}
\newtheorem{corollary}[definition]{Corollary}
\newtheorem{interpretation}[definition]{Interpretation}
\newtheorem{empirical}[definition]{Empirical statement}
\newtheorem{established}[definition]{Established physical result}
\newtheorem{newresult}[definition]{New result claimed by this work}
\newtheorem{counterexample}[definition]{Counterexample}
\newtheorem{example}[definition]{Example}

\newcommand{\R}{\mathbb{R}}
\newcommand{\Lie}{\mathcal{L}}

\title{Layered Composition, Reciprocity, and Protocol-Relative Realization}
\author{PRINCIPIA PHYSICA Synthesis}
\date{15 August 2026}

\begin{document}
\maketitle

\begin{abstract}
This synthesis separates mathematical admissibility from physical realization.
It defines a layered realized theory object, proves a reciprocity trichotomy for
a linear interconnection, and proves that closed-system isomorphism need not
preserve a fixed empirical protocol. It also corrects an earlier claim about
limits and composition: for the oscillator family, the zero-frequency limit
commutes with composition when the interaction parameter is retained. A
different trajectory arises only when that interface datum is discarded.
Lean and Haskell companions formalize finite logical and computational shadows
of these claims; neither artifact formalizes the differential geometry in full.
\end{abstract}

\section{Epistemic discipline}

\begin{assumption}[Finite-dimensional scope, S-A1]
All state spaces considered below are finite-dimensional smooth manifolds.
Claims about field theory, general relativity, and quantum theory are outside
the proved scope.
\end{assumption}

\begin{axiom}[Synthesis discipline, S-A2]
A synthesis may translate inherited claims and prove consequences from jointly
satisfiable hypotheses. It may not promote an inherited epistemic label, treat
agreement as evidence, or identify mathematical admissibility with empirical
realization.
\end{axiom}

\begin{interpretation}[Founding thesis, S-I1]
The thesis that physics studies empirically realized compositional mathematical
structures is treated as a methodological proposal. The mathematical results
below test distinctions required by that proposal; they do not establish its
exhaustiveness or empirical truth.
\end{interpretation}

\section{Layered realized systems}

\begin{definition}[Realized theory object, S-D1]
A realized theory object is a tuple
\[
  \mathcal R=(K,D,I,E,P,A),
\]
where $K$ contains states, observables, and declared geometry; $D$ contains a
generator and its flow domain; $I$ contains ports and interaction data; $E$
contains preparations, measurement kernels, calibration, and uncertainty;
$P$ is a protocol; and $A$ is a predeclared adequacy rule. Mathematical
admissibility concerns the appropriate mathematical components. Physical
realization additionally requires empirical input in $E$, $P$, and $A$.
\end{definition}

\begin{proposition}[Vocabulary translation is not top-level injective, S-P1]
The assignment in Table~\ref{tab:translation} is a function on the listed
source notions, but it is not injective into the set
$\{K,D,I,E,P,A\}$. In particular, kinematical product and observable embedding
both map to $K$, while preparation distributions and measurement kernels both
map to $E$.
\end{proposition}

\begin{proof}
Every row in Table~\ref{tab:translation} has one source and one target, so the
displayed assignment is a function. The first two rows have the same target
$K$, and the last two rows have the same target $E$. Hence the assignment is
not injective into the six top-level components. The parenthetical qualifiers
identify distinct substructures within those components; they do not create
additional top-level tuple entries.
\end{proof}

\begin{table}[h]
\centering
\begin{tabular}{lll}
\toprule
Source notion & Target & Internal field \\
\midrule
Kinematical product & $K$ & state space \\
Observable embedding & $K$ & observable algebra \\
Dynamical factorization & $D$ & generator and flow \\
Port observable & $I$ & distinguished observable \\
Preparation distribution & $E$ & preparation field \\
Measurement kernel & $E$ & measurement field \\
Adequacy rule & $A$ & acceptance region \\
\bottomrule
\end{tabular}
\caption{A non-injective translation into the layered carrier.}
\label{tab:translation}
\end{table}

\section{Reciprocity and factorization}

Let $X=\R^{2n_1}\times\R^{2n_2}$ have canonical product symplectic form
$\omega_\oplus$. For matrices $C_{12}$ and $C_{21}$, define
\[
 Y=(0,-C_{12}q_2,0,-C_{21}q_1).
\]

\begin{lemma}[Contraction calculation, S-L1]
For the stated field,
\[
 \iota_Y\omega_\oplus
 =-(C_{12}q_2)\cdot dq_1-(C_{21}q_1)\cdot dq_2,
\]
and
\[
 d(\iota_Y\omega_\oplus)
 =\sum_{a,b}\bigl((C_{21})_{ba}-(C_{12})_{ab}\bigr)
   dq_{1a}\wedge dq_{2b}.
\]
\end{lemma}

\begin{proof}
For a canonical pair, contraction of $(\dot q,\dot p)$ with
$dq\wedge dp$ is $\dot q\,dp-\dot p\,dq$. Substitution of the two momentum
components of $Y$ gives the first display. Differentiating each coefficient and
reordering $dq_{2b}\wedge dq_{1a}=-dq_{1a}\wedge dq_{2b}$ gives the second.
\end{proof}

\begin{newresult}[Linear composition trichotomy, S-N1]
Relative to $\omega_\oplus$, the stated family has three disjoint classes:
(i) $C_{12}=C_{21}=0$, for which the dynamics factorize; (ii)
$C_{21}=C_{12}^{T}\ne0$, for which the interconnection is Hamiltonian but does
not factorize; and (iii) $C_{21}\ne C_{12}^{T}$, for which the interconnection
does not preserve $\omega_\oplus$.
\end{newresult}

\begin{proof}
By Cartan's formula and $d\omega_\oplus=0$,
$\Lie_Y\omega_\oplus=d(\iota_Y\omega_\oplus)$. Lemma S-L1 makes this zero
exactly when $C_{21}=C_{12}^{T}$. If both matrices vanish, $Y=0$ and an
additive factor generator remains. If the reciprocal matrices are nonzero,
$Y$ has a component depending on the other factor, so the vector field is not
a product field. If reciprocity fails, the displayed Lie derivative is
nonzero. These conditions are mutually exclusive and exhaustive.
\end{proof}

\begin{corollary}[Hamiltonian does not imply factorized, S-C1]
The decoupled class is a proper subset of the Hamiltonian class whenever a
nonzero matrix $C$ is available: take $C_{12}=C$ and $C_{21}=C^T$.
\end{corollary}

\begin{example}[Normal modes expose the middle class]
For one degree of freedom in each factor, take
$C_{12}=C_{21}=c\ne0$ and add identical oscillator Hamiltonians.  The equations
for $q_\pm=(q_1\pm q_2)/\sqrt2$ decouple into normal modes with squared
frequencies $\omega^2\pm c/m$.  Thus the total system has a Hamiltonian flow,
yet neither original factor evolves without reference to the other.  The
example witnesses class (ii) of S-N1 without relying only on the abstract
matrix criterion.
\end{example}

\begin{counterexample}[Volume preservation is not reciprocity]
The divergence of $Y$ is zero for every pair $C_{12},C_{21}$ because its
Jacobian has zero diagonal.  Taking $C_{12}=0$ and $C_{21}\ne0$ therefore
gives a volume-preserving interconnection in class (iii).  Preservation of
phase volume alone cannot replace preservation of $\omega_\oplus$.
\end{counterexample}

\section{Closed isomorphism and empirical protocol}

\begin{newresult}[Protocol distinction survives closed isomorphism, S-N2]
Let
\[
 H_m(q,p)=\frac{p^2}{2m}+\frac{m\omega_0^2q^2}{2}.
\]
For positive $m,m'$, the map
$F(q,p)=(\alpha q,p/\alpha)$ with $\alpha=\sqrt{m/m'}$ is symplectic and
satisfies $H_{m'}\circ F=H_m$. Nevertheless, a fixed protocol preparing
$(q,p)=(0,1)$ and reading the unscaled coordinate $q(t)$ distinguishes the
masses whenever the two predicted readings differ.
\end{newresult}

\begin{proof}
The pullback of $dq\wedge dp$ under $F$ is
$d(\alpha q)\wedge d(p/\alpha)=dq\wedge dp$. Direct substitution gives
$H_{m'}(F(q,p))=H_m(q,p)$. The prepared solution is
$q_m(t)=\sin(\omega_0t)/(m\omega_0)$. Thus, for example, $m=1$, $m'=4$,
$\omega_0=1$, and $t=1$ give readings $\sin 1$ and $\sin 1/4$, which are
unequal. The mathematical isomorphism does not transport the fixed calibration
used by the protocol.
\end{proof}

\begin{interpretation}[Calibration transport, S-I2]
Transporting the port calibration with $F$ removes this particular protocol
difference. Keeping the calibration fixed does not. This is an interpretation
of S-N2, not an empirical assertion that either oscillator is physically
realized.
\end{interpretation}

\begin{proposition}[Transported protocol restores equivalence, S-P2]
Let a protocol for $H_m$ have preparation measure $\pi$, evolution $\phi_t$,
and measurement kernel $K$.  Define the protocol for $H_{m'}$ by
$\pi'=F_*\pi$, $\phi'_t=F\phi_tF^{-1}$, and
$K'(x,B)=K(F^{-1}x,B)$.  Then the two predicted outcome measures agree for
every admitted time and event $B$.
\end{proposition}
\begin{proof}
The primed prediction is
\[
 \int K'(\phi'_t(x),B)\,d\pi'(x)
 =\int K(F^{-1}F\phi_t(y),B)\,d\pi(y)
 =\int K(\phi_t(y),B)\,d\pi(y),
\]
where the first equality uses the pushforward change of variables.  This is
the unprimed prediction.  S-N2 differs because it intentionally fixes rather
than transports the preparation and coordinate calibration.
\end{proof}

\begin{counterexample}[Partial transport is insufficient]
Transport the Hamiltonian and preparation by $F$ but retain the original
coordinate-reading kernel.  The evolved mathematical states correspond, but
the unscaled readout again returns $q_m(t)$ versus $q_{m'}(t)$.  Dynamical
intertwining alone does not establish protocol equivalence.
\end{counterexample}

\section{Limits and interface retention}

\begin{newresult}[Fixed-interface commutation, S-N3]
Let $\Sigma^{(i)}_\omega$ be identical oscillators of mass $m$ and frequency
$\omega$, and let $\otimes_\kappa$ add the interaction $\kappa q_1q_2$. Then
\[
 \lim_{\omega\to0}
 (\Sigma^{(1)}_\omega\otimes_\kappa\Sigma^{(2)}_\omega)
 =
 (\lim_{\omega\to0}\Sigma^{(1)}_\omega)
 \otimes_\kappa
 (\lim_{\omega\to0}\Sigma^{(2)}_\omega).
\]
If $\otimes_0$ is used on the right instead, its Hamiltonian differs from the
fixed-interface limit by $\kappa q_1q_2$.
\end{newresult}

\begin{proof}
Before taking the limit, the composite Hamiltonian is
\[
 H_{\omega,\kappa}
 =\frac{p_1^2+p_2^2}{2m}
  +\frac{m\omega^2(q_1^2+q_2^2)}2+\kappa q_1q_2.
\]
Taking $\omega\to0$ gives
$H_{0,\kappa}=(p_1^2+p_2^2)/(2m)+\kappa q_1q_2$. Taking the isolated limits
first gives $p_i^2/(2m)$, and applying the same $\otimes_\kappa$ adds the same
interaction term. Both orders therefore yield $H_{0,\kappa}$. Applying
$\otimes_0$ after the factor limits instead yields $H_{0,0}$, whose difference
from $H_{0,\kappa}$ is $\kappa q_1q_2$.
\end{proof}

\begin{counterexample}[Counterexample to the rejected non-commutation claim]
Retain $\kappa$ in both composition operations. The two Hamiltonians are then
identical. Thus a comparison using $\otimes_\kappa$ before the limit and
$\otimes_0$ afterward does not establish non-commutation of a fixed operation;
it establishes that changing the interface data changes the system.
\end{counterexample}

\begin{proposition}[Vector-field form of fixed-interface commutation, S-P3]
For the Hamiltonians in S-N3, the Hamiltonian vector fields converge
coefficientwise, uniformly on compact subsets of phase space, to the vector
field of $H_{0,\kappa}$.  Taking isolated limits and then applying the same
$\otimes_\kappa$ produces that identical limiting field.
\end{proposition}
\begin{proof}
Hamilton's equations are
\[
 \dot q_i=p_i/m,\qquad
 \dot p_1=-m\omega^2q_1-\kappa q_2,\qquad
 \dot p_2=-m\omega^2q_2-\kappa q_1.
\]
On a compact set with $|q_i|\le R$, the difference from the $\omega=0$ field
has norm at most $\sqrt2m\omega^2R$.  This tends uniformly to zero.  Applying
the retained interface after the isolated limits writes down the same last
two terms, so it gives exactly the same vector field.
\end{proof}

\begin{interpretation}[What the estimate does not prove]
Uniform convergence of vector fields on compact phase-space sets supports
finite-time trajectory comparison only together with an ODE stability
argument and a common compact region containing the trajectories.  It does
not establish uniform agreement for all initial data or all times.  The
bounded-domain warning inherited from EL therefore remains relevant even
though the algebraic commutation statement is exact.
\end{interpretation}

\begin{proposition}[Finite-time trajectory estimate, S-P5]
Fix a compact convex phase-space region $B$ and $T>0$.  Suppose the
$\omega$ and zero-frequency solutions with common initial data remain in $B$
for $0\le t\le T$.  Then, for a constant $L$ bounding the Lipschitz constant
of the limiting linear vector field on $B$,
\[
 \|z_\omega(t)-z_0(t)\|
 \le \frac{\sqrt2m\omega^2R}{L}(e^{Lt}-1)
\]
when $L>0$, with the quotient replaced by $t$ when $L=0$.
\end{proposition}
\begin{proof}
Write the vector fields as $f_\omega$ and $f_0$.  The calculation in S-P3
gives $\|f_\omega(z)-f_0(z)\|\le\varepsilon_\omega:=\sqrt2m\omega^2R$ on
$B$.  The integral equations imply
\[
 \|z_\omega(t)-z_0(t)\|
 \le\int_0^t\bigl(\varepsilon_\omega
       +L\|z_\omega(s)-z_0(s)\|\bigr)\,ds.
\]
Gr\"onwall's inequality yields the displayed estimate.  The explicit
compact-region hypothesis prevents the finite-time statement from being
mistaken for a global uniform limit.
\end{proof}

\begin{counterexample}[Long times evade the estimate]
Even for a single uncoupled oscillator,
$q_\omega(t)=q_0\cos(\omega t)+(v_0/\omega)\sin(\omega t)$ approaches the free
trajectory on bounded time intervals but not uniformly for all time.  Taking
$t=2\pi/\omega$ with $v_0\ne0$ produces a discrepancy
$2\pi|v_0|/\omega$.  Exact commutation of the Hamiltonian formula in S-N3
therefore does not erase the domain qualifications attached to convergence of
solutions.
\end{counterexample}

\subsection{Protocol quotient and composition}

For a fixed protocol, observational equivalence identifies realized models
having identical families of predicted outcome measures.  Composition can
descend to this quotient only if replacing either factor by an equivalent
representative leaves every composite prediction unchanged.  This is a
congruence requirement, stronger than equivalence of isolated predictions.

\begin{definition}[Protocol congruence, S-D3]
An observational equivalence $\sim_P$ is a congruence for an
interface-indexed operation $\otimes_\kappa$ when
$M\sim_PM'$ and $N\sim_PN'$ imply
$M\otimes_\kappa N\sim_{P_\kappa}M'\otimes_\kappa N'$ for the declared
composite protocol $P_\kappa$.
\end{definition}

\begin{counterexample}[Isolated equivalence need not be a congruence]
Let an isolated protocol read only positions at time zero.  Two oscillator
models with identical prepared position distributions but different momenta
are equivalent under that impoverished protocol.  Couple each to a second
system and later read position.  Hamiltonian evolution can convert the hidden
momentum difference into a position difference, so the composites are
distinguishable.  Equivalence under the isolated protocol is therefore not
automatically substitutive under richer composite interventions.
\end{counterexample}

\begin{interpretation}[Required interface information]
The counterexample does not show that quotient composition is impossible.  It
shows that a protocol must expose enough port behavior, or equivalence must be
defined against a sufficiently rich family of interventions, before
congruence can be expected.  This motivates the interface datum in S-D1 and
the qualification ``suitable'' in S-C3.
\end{interpretation}

\begin{empirical}[Computational cross-check, S-E2]
For $m=\kappa=1$ and initial data $(0,0,1,0)$, the retained-interface system
has $q_1(t)=\tfrac12(\cos t-\cosh t)$, while the discarded-interface system
has $q_1(t)=0$. The accompanying Haskell program evaluates both at $t=1$ by
fourth-order Runge--Kutta integration. This finite-precision calculation is a
software-process record, not a proof and not an observation of a physical
system.
\end{empirical}

\section{Formal scope and dependencies}

\begin{definition}[Dependency ledger, S-D2]
A dependency ledger records, for each inherited claim used by the synthesis,
its source identifier, epistemic type, hypotheses, and the exact conclusion
imported.  Recording a dependency does not re-prove it; it makes the boundary
between inherited and newly proved content inspectable.
\end{definition}

\begin{proposition}[No empirical promotion by conjunction, S-P4]
Conjoining finitely many mathematical theorems and interpretations cannot by
itself entail an empirical realization statement under S-A2 unless an
empirical premise connecting the model to a procedure or observation is also
supplied.
\end{proposition}
\begin{proof}
S-A2 forbids identifying mathematical admissibility with realization.  The
conclusions of the mathematical results constrain $K,D,$ and $I$, while a
realization assertion additionally concerns $E,P,$ and $A$ by S-D1.  Without
a premise about those components, the required connection is absent.  This is
a typed consequence of the stipulated synthesis discipline, not a theorem of
untyped propositional logic.
\end{proof}

\begin{counterexample}[Agreement of artifacts is not independent evidence]
If a Lean registry and a Haskell test are both transcribed from the same false
paper premise, they may agree with one another while inheriting the same
error.  Their agreement checks consistency of selected encodings, not the
truth of an empirical premise or the correctness of omitted differential
geometry.  S-E3 therefore records their limited scope explicitly.
\end{counterexample}

\begin{empirical}[Tool artifact record, S-E3]
The accompanying \texttt{Proofs.lean} uses Lean core to check a finite claim
registry and algebraic witnesses. The accompanying \texttt{Core.hs} provides
finite numerical checks. These files do not formalize smooth manifolds,
differential forms, ODE existence, or physical realization.
\end{empirical}

\begin{conjecture}[Interface-indexed composition, S-C3]
There may exist a category of realized theory objects in which composition is
indexed by explicit interface data and is compatible with protocol quotients
and suitable limits. No such category, coherence theorem, or general
continuity theorem is constructed here.
\end{conjecture}

\section{Bibliographic notes}

The dependency identifiers CS, HR, and EL refer respectively to the supplied
papers \emph{Compositional Systems in Finite-Dimensional Classical Mechanics},
\emph{Hamiltonian Reconstruction, Reciprocity, and Typed Physical Structure},
and \emph{Empirical Realization, Equivalence, Limits, and Falsification}.
CS-P4 is used only for the statement that isolated factors do not determine an
interaction. EL-P3 and EL-P4 are cited only for their bounded-domain limit
warnings. The 15 August 2026 synthesis review supplied the counterexample that
forced the corrected formulation of S-N3 and identified the non-injectivity
error in the former S-P1.

\section{Conclusion}

The synthesis establishes three scoped mathematical points. Reciprocal
coupling can be Hamiltonian without being factorized. Closed Hamiltonian
isomorphism does not determine equivalence under a fixed empirical protocol.
For the oscillator family, a limit commutes with a fixed interaction-indexed
composition, whereas discarding the interaction produces a different system.
None of these mathematical results supplies evidence that a modeled structure
is empirically realized.

\end{document}
